\section*{Chapter \ref{ch:b1:predominant-reaction} --- Acid--Base Equilibria and the Predominant-Reaction Method}

\begin{solution}{exo:b1:predominant-reaction:1}
Conjugate bases: \ce{HSO4-}, \ce{CO3^{2-}}, \ce{NH3}, \ce{OH-}. Conjugate
acids: \ce{NH4+}, \ce{H2PO4-}, \ce{H2O}, \ce{H3O+}. \omterm{def:b1:predominant-reaction:ampholyte}{Ampholytes}: \ce{HCO3-},
\ce{H2O}, \ce{HPO4^{2-}} (and \ce{HSO4-}, base of sulfuric acid and acid of
its second couple).
\end{solution}

\begin{solution}{exo:b1:predominant-reaction:2}
\ce{HF} ($K_a = \num{6.9e-4}$) $>$ \ce{CH3COOH} (\num{1.7e-5}) $>$ \ce{HClO}
(\num{2.8e-8}) $>$ \ce{NH4+} (\num{5.6e-10}).
\end{solution}

\begin{solution}{exo:b1:predominant-reaction:3}
\omterm{def:b1:predominant-reaction:strong-acid}{Strong acid}: $\mathrm{pH} = 2.00$. \omterm{def:b1:predominant-reaction:strong-acid}{Strong base}: $[\ce{OH-}] = 10^{-2}$,
$\mathrm{pH} = 14.00 - 2.00 = 12.00$.
\end{solution}

\begin{solution}{exo:b1:predominant-reaction:4}
\ce{CH3COOH} below pH 4.76, \ce{CH3COO-} above. At pH 3 the acid form
predominates (98\,\%); in blood (pH 7.4) the ethanoate ion. At pH 3.76 the
acid fraction is $1/(1 + 10^{-1}) = 0.91$.
\end{solution}

\begin{solution}{exo:b1:predominant-reaction:5}
\omterm{def:b1:predominant-reaction:predominant-reaction}{Predominant reaction} \ce{NH3 + H2O <=> NH4+ + OH-}, $K = 10^{-(14.00 -
9.25)} = 10^{-4.75}$. If little \ce{NH3} reacts: $[\ce{OH-}] = \sqrt{K C} =
\qty{1.3e-3}{mol/L}$, $\mathrm{pH} = 11.13$. Check: $[\ce{NH4+}] = 1.3\,\%$
of $C$, and $\mathrm{pH} > \mathrm{p}K_a + 1$: valid.
\end{solution}

\begin{solution}{exo:b1:predominant-reaction:6}
$\mathrm{pH} = \frac12(4.76 + 2.00) = 3.38$ (the exact value is 3.39);
$\alpha = 10^{-3.38}/0.010 = 0.042$: 4\,\% dissociated, below 10\,\%, and
$\mathrm{pH} < \mathrm{p}K_a - 1$: valid.
\end{solution}

\begin{solution}{exo:b1:predominant-reaction:7}
\ce{HF + OH- -> F- + H2O}, $K = 10^{14.00 - 3.16} = 10^{10.84}$: total. The
\omterm{def:b1:predominant-reaction:predominant-reaction}{equivalent solution} holds \qty{0.050}{mol/L} each of \ce{HF} and \ce{F-}, a
buffer: $\mathrm{pH} = \mathrm{p}K_a = 3.16$.
\end{solution}

\begin{solution}{exo:b1:predominant-reaction:8}
$\mathrm{pH} = 4.76$. The hydrochloric acid converts \qty{0.010}{mol} of
ethanoate into acid: $\mathrm{pH} = 4.76 + \log(0.040/0.060) = 4.58$, a
change of 0.18. In pure water the pH would fall from 7.00 to 2.00.
\end{solution}

\begin{solution}{exo:b1:predominant-reaction:9}
In water both acids react totally to give \ce{H3O+}: the solvent levels
them. Ethanoic acid is a much weaker base than water, so in it the two
acids react only partly, to different extents, and their strengths can be
compared. The strongest acid that can exist in water is \ce{H3O+}.
\end{solution}

\begin{solution}{exo:b1:predominant-reaction:10}
\ce{CO3^{2-} + H2O <=> HCO3- + OH-}, $K = 10^{-(14.00 - 10.33)} = 10^{-3.67}$;
$[\ce{OH-}] = \sqrt{10^{-3.67} \times 0.10} = \qty{4.6e-3}{mol/L}$ (4.6\,\%
of $C$), $\mathrm{pH} = 11.67$. The second basicity, \ce{HCO3- + H2O <=>
CO2 + H2O + OH-}, has $K = 10^{-7.63}$: it gives $[\ce{CO2}] \approx
10^{-7.6}$, negligible.
\end{solution}

\begin{solution}{exo:b1:predominant-reaction:11}
$[\mathrm{A}] = [\ce{H3O+}] = \alpha C$ and $[\mathrm{HA}] = (1 - \alpha)C$,
so $K_a = \alpha^2C/(1 - \alpha)$. With $C = 10^{-3}$: $\alpha^2 +
0.0174\,\alpha - 0.0174 = 0$, $\alpha = 0.12$. The acid is 12\,\%
dissociated: the hypothesis of a little-dissociated acid ($\alpha <
10\,\%$) fails, and the quadratic must be solved.
\end{solution}

\begin{solution}{exo:b1:predominant-reaction:12}
Hydrochloric acid fixes $[\ce{H3O+}] = \qty{0.010}{mol/L}$: pH 2.00. Then
\[
  [\ce{CH3COO-}] = \frac{K_a[\ce{CH3COOH}]}{[\ce{H3O+}]} = \frac{1.74 \times 10^{-5}
  \times 0.10}{0.010} = \qty{1.7e-4}{mol/L},
\]
negligible beside 0.010: the
\ce{H3O+} from the \omterm{def:b1:predominant-reaction:strong-acid}{strong acid} pushes the equilibrium of the \omterm{def:b1:predominant-reaction:strong-acid}{weak acid}
back.
\end{solution}

\begin{solution}{pb:b1:predominant-reaction:1}
\textbf{1.} \ce{H3PO4}/\ce{H2PO4-}, \ce{H2PO4-}/\ce{HPO4^{2-}},
\ce{HPO4^{2-}}/\ce{PO4^{3-}}; for instance \ce{H3PO4 <=> H2PO4- + H+}.
\textbf{2.} Boundaries at 2.15, 7.21 and 12.34.
\textbf{3.} pH 1: \ce{H3PO4}; 5: \ce{H2PO4-}; 9: \ce{HPO4^{2-}}; 13:
\ce{PO4^{3-}}.
\textbf{4.} \ce{H2PO4-} and \ce{HPO4^{2-}}.
\textbf{5.} Each proton leaves an ion of higher negative charge, which
holds the next proton more strongly.
\textbf{6.} \ce{H3PO4 + H2O <=> H2PO4- + H3O+}; the formula
$\frac12(\mathrm{p}K_a + \mathrm{p}C) = 1.58$ violates $\mathrm{pH} <
\mathrm{p}K_a - 1 = 1.15$: about a fifth of the acid reacts.
\textbf{7.} $x^2 + K_{a1}x - K_{a1}C = 0$ with $K_{a1} = 10^{-2.15}$: $x =
\qty{0.0233}{mol/L}$, $\mathrm{pH} = 1.63$.
\textbf{8.} \ce{2H2PO4- <=> H3PO4 + HPO4^{2-}}, $K = 10^{2.15 - 7.21} =
10^{-5.06}$.
\textbf{9.} $\mathrm{pH} = \frac12(2.15 + 7.21) = 4.68$.
\textbf{10.} $\mathrm{pH} = \frac12(7.21 + 12.34) = 9.78$.
\textbf{11.} None: 1.63, 4.68 and 9.78; a mixture of two of them is
needed.
\textbf{12.} \ce{H3PO4 + OH- -> H2PO4- + H2O}, $K = 10^{14.00 - 2.15} =
10^{11.85}$: \omterm{def:b1:extent-q-and-k:quantitative}{quantitative}.
\textbf{13.} \qty{0.100}{mol/L} of \ce{H2PO4-}: pH 4.68.
\textbf{14.} The next \qty{0.050}{mol} of hydroxide converts half of it:
\qty{0.050}{mol/L} of \ce{H2PO4-} and of \ce{HPO4^{2-}}, pH 7.21.
\textbf{15.} \qty{0.100}{mol/L} of \ce{HPO4^{2-}}: pH 9.78.
\textbf{16.} The pH rises from 1.63, jumps around $n = 0.100$ (to 4.68),
rises slowly through 7.21 at $n = 0.150$, jumps around $n = 0.200$ (to
9.78), passes 12.34 near $n = 0.250$ and reaches about 12.6 at $n = 0.300$ (\qty{0.100}{mol/L} of \ce{PO4^3-}: $x^2/(0.100 - x) = 10^{-1.66}$, $x = 0.037$, pH 12.57).
\textbf{17.} It contains equal amounts of the acid and the base of a couple
of $\mathrm{p}K_a$ 7.21 (\cref{prop:b1:predominant-reaction:buffer}).
\textbf{18.} $10^{7.20 - 7.21} = 0.977$.
\textbf{19.} $n(\ce{H2PO4-}) = 0.100 - x$, $n(\ce{HPO4^{2-}}) = x$.
\textbf{20.} $x/(0.100 - x) = 0.977$, $x = \qty{0.0494}{mol}$.
\textbf{21.} The acid converts \qty{0.0010}{mol} of \ce{HPO4^{2-}}:
$\mathrm{pH} = 7.21 + \log(0.0484/0.0516) = 7.18$, a change of $-0.02$. In
pure water: from 7.00 to 3.00.
\textbf{22.} No: the pH depends on the ratio of the amounts, which dilution
does not change (as long as the activities stay close to the
concentrations).
\textbf{23.} $0.100 + 0.0494 = \qty{0.149}{mol}$ of hydroxide: \textbf{\qty{149}{mL}}
of the \qty{1.00}{mol/L} solution, made up to \qty{1.00}{L}.
\end{solution}
